\begin{abstract}
Formula~1 team radio pairs short, noisy, jargon-heavy speech with an independent record of what the car was doing. We build an end-to-end, pre-registered pipeline over 14{,}370 broadcast clips (2018--2025) joined to telemetry for 148 races, and study claim-type classification, telemetry consistency, vocal arousal, and arousal's effect on subsequent laps. Transfer-learning choices follow each model's data budget: full fine-tuning of DeBERTa-v3 with a shared multi-task trunk, partial fine-tuning of WavLM with a learned layer weighting, and prompt-only adaptation of Whisper verified by word error rate. Evaluation holds out whole races and whole drivers and repeats every telemetry test on a reverse-time placebo window. Text identifies claim type well above a context-only baseline (macro AUC-PR 0.43 vs.\ 0.32) and generalises to unseen drivers. With approximately 250 labels, the fine-tuned transformer fails to outperform TF-IDF, while a small zero-shot LLM utilizing the category definitions surpasses both (0.63). Partial fine-tuning lifts audio arousal from 0.44 to 0.55. Speaker attribution scores 0.67 against a target of 0.85, and prompt biasing decreases WER by only about 3\%. Telemetry supports no claim type beyond its base rate: wherever a signal is detected, the placebo indicates it precedes the message, consistent with broadcast selection. Arousal shows no detectable effect on subsequent pace.
\end{abstract}
